% Options for packages loaded elsewhere
% Options for packages loaded elsewhere
\PassOptionsToPackage{unicode}{hyperref}
\PassOptionsToPackage{hyphens}{url}
%
\documentclass[
  ngerman,
  10pt,
  letterpaper,
  twocolumn]{book}
\usepackage{xcolor}
\usepackage[top=2.5cm,bottom=2.5cm,left=2cm,right=2cm]{geometry}
\usepackage{amsmath,amssymb}
\setcounter{secnumdepth}{5}
\usepackage{iftex}
\ifPDFTeX
  \usepackage[T1]{fontenc}
  \usepackage[utf8]{inputenc}
  \usepackage{textcomp} % provide euro and other symbols
\else % if luatex or xetex
  \usepackage{unicode-math} % this also loads fontspec
  \defaultfontfeatures{Scale=MatchLowercase}
  \defaultfontfeatures[\rmfamily]{Ligatures=TeX,Scale=1}
\fi
\usepackage{lmodern}
\ifPDFTeX\else
  % xetex/luatex font selection
\fi
% Use upquote if available, for straight quotes in verbatim environments
\IfFileExists{upquote.sty}{\usepackage{upquote}}{}
\IfFileExists{microtype.sty}{% use microtype if available
  \usepackage[]{microtype}
  \UseMicrotypeSet[protrusion]{basicmath} % disable protrusion for tt fonts
}{}
\usepackage{setspace}
\makeatletter
\@ifundefined{KOMAClassName}{% if non-KOMA class
  \IfFileExists{parskip.sty}{%
    \usepackage{parskip}
  }{% else
    \setlength{\parindent}{0pt}
    \setlength{\parskip}{6pt plus 2pt minus 1pt}}
}{% if KOMA class
  \KOMAoptions{parskip=half}}
\makeatother
% Make \paragraph and \subparagraph free-standing
\makeatletter
\ifx\paragraph\undefined\else
  \let\oldparagraph\paragraph
  \renewcommand{\paragraph}{
    \@ifstar
      \xxxParagraphStar
      \xxxParagraphNoStar
  }
  \newcommand{\xxxParagraphStar}[1]{\oldparagraph*{#1}\mbox{}}
  \newcommand{\xxxParagraphNoStar}[1]{\oldparagraph{#1}\mbox{}}
\fi
\ifx\subparagraph\undefined\else
  \let\oldsubparagraph\subparagraph
  \renewcommand{\subparagraph}{
    \@ifstar
      \xxxSubParagraphStar
      \xxxSubParagraphNoStar
  }
  \newcommand{\xxxSubParagraphStar}[1]{\oldsubparagraph*{#1}\mbox{}}
  \newcommand{\xxxSubParagraphNoStar}[1]{\oldsubparagraph{#1}\mbox{}}
\fi
\makeatother

\usepackage{color}
\usepackage{fancyvrb}
\newcommand{\VerbBar}{|}
\newcommand{\VERB}{\Verb[commandchars=\\\{\}]}
\DefineVerbatimEnvironment{Highlighting}{Verbatim}{commandchars=\\\{\}}
% Add ',fontsize=\small' for more characters per line
\usepackage{framed}
\definecolor{shadecolor}{RGB}{241,243,245}
\newenvironment{Shaded}{\begin{snugshade}}{\end{snugshade}}
\newcommand{\AlertTok}[1]{\textcolor[rgb]{0.68,0.00,0.00}{#1}}
\newcommand{\AnnotationTok}[1]{\textcolor[rgb]{0.37,0.37,0.37}{#1}}
\newcommand{\AttributeTok}[1]{\textcolor[rgb]{0.40,0.45,0.13}{#1}}
\newcommand{\BaseNTok}[1]{\textcolor[rgb]{0.68,0.00,0.00}{#1}}
\newcommand{\BuiltInTok}[1]{\textcolor[rgb]{0.00,0.23,0.31}{#1}}
\newcommand{\CharTok}[1]{\textcolor[rgb]{0.13,0.47,0.30}{#1}}
\newcommand{\CommentTok}[1]{\textcolor[rgb]{0.37,0.37,0.37}{#1}}
\newcommand{\CommentVarTok}[1]{\textcolor[rgb]{0.37,0.37,0.37}{\textit{#1}}}
\newcommand{\ConstantTok}[1]{\textcolor[rgb]{0.56,0.35,0.01}{#1}}
\newcommand{\ControlFlowTok}[1]{\textcolor[rgb]{0.00,0.23,0.31}{\textbf{#1}}}
\newcommand{\DataTypeTok}[1]{\textcolor[rgb]{0.68,0.00,0.00}{#1}}
\newcommand{\DecValTok}[1]{\textcolor[rgb]{0.68,0.00,0.00}{#1}}
\newcommand{\DocumentationTok}[1]{\textcolor[rgb]{0.37,0.37,0.37}{\textit{#1}}}
\newcommand{\ErrorTok}[1]{\textcolor[rgb]{0.68,0.00,0.00}{#1}}
\newcommand{\ExtensionTok}[1]{\textcolor[rgb]{0.00,0.23,0.31}{#1}}
\newcommand{\FloatTok}[1]{\textcolor[rgb]{0.68,0.00,0.00}{#1}}
\newcommand{\FunctionTok}[1]{\textcolor[rgb]{0.28,0.35,0.67}{#1}}
\newcommand{\ImportTok}[1]{\textcolor[rgb]{0.00,0.46,0.62}{#1}}
\newcommand{\InformationTok}[1]{\textcolor[rgb]{0.37,0.37,0.37}{#1}}
\newcommand{\KeywordTok}[1]{\textcolor[rgb]{0.00,0.23,0.31}{\textbf{#1}}}
\newcommand{\NormalTok}[1]{\textcolor[rgb]{0.00,0.23,0.31}{#1}}
\newcommand{\OperatorTok}[1]{\textcolor[rgb]{0.37,0.37,0.37}{#1}}
\newcommand{\OtherTok}[1]{\textcolor[rgb]{0.00,0.23,0.31}{#1}}
\newcommand{\PreprocessorTok}[1]{\textcolor[rgb]{0.68,0.00,0.00}{#1}}
\newcommand{\RegionMarkerTok}[1]{\textcolor[rgb]{0.00,0.23,0.31}{#1}}
\newcommand{\SpecialCharTok}[1]{\textcolor[rgb]{0.37,0.37,0.37}{#1}}
\newcommand{\SpecialStringTok}[1]{\textcolor[rgb]{0.13,0.47,0.30}{#1}}
\newcommand{\StringTok}[1]{\textcolor[rgb]{0.13,0.47,0.30}{#1}}
\newcommand{\VariableTok}[1]{\textcolor[rgb]{0.07,0.07,0.07}{#1}}
\newcommand{\VerbatimStringTok}[1]{\textcolor[rgb]{0.13,0.47,0.30}{#1}}
\newcommand{\WarningTok}[1]{\textcolor[rgb]{0.37,0.37,0.37}{\textit{#1}}}

\usepackage{longtable,booktabs,array}
\usepackage{calc} % for calculating minipage widths
% Correct order of tables after \paragraph or \subparagraph
\usepackage{etoolbox}
\makeatletter
\patchcmd\longtable{\par}{\if@noskipsec\mbox{}\fi\par}{}{}
\makeatother
% Allow footnotes in longtable head/foot
\IfFileExists{footnotehyper.sty}{\usepackage{footnotehyper}}{\usepackage{footnote}}
\makesavenoteenv{longtable}
\usepackage{graphicx}
\makeatletter
\newsavebox\pandoc@box
\newcommand*\pandocbounded[1]{% scales image to fit in text height/width
  \sbox\pandoc@box{#1}%
  \Gscale@div\@tempa{\textheight}{\dimexpr\ht\pandoc@box+\dp\pandoc@box\relax}%
  \Gscale@div\@tempb{\linewidth}{\wd\pandoc@box}%
  \ifdim\@tempb\p@<\@tempa\p@\let\@tempa\@tempb\fi% select the smaller of both
  \ifdim\@tempa\p@<\p@\scalebox{\@tempa}{\usebox\pandoc@box}%
  \else\usebox{\pandoc@box}%
  \fi%
}
% Set default figure placement to htbp
\def\fps@figure{htbp}
\makeatother



\ifLuaTeX
\usepackage[bidi=basic]{babel}
\else
\usepackage[bidi=default]{babel}
\fi
% get rid of language-specific shorthands (see #6817):
\let\LanguageShortHands\languageshorthands
\def\languageshorthands#1{}
\ifLuaTeX
  \usepackage[german]{selnolig} % disable illegal ligatures
\fi


\setlength{\emergencystretch}{3em} % prevent overfull lines

\providecommand{\tightlist}{%
  \setlength{\itemsep}{0pt}\setlength{\parskip}{0pt}}



 
\usepackage[style=authoryear,backend=biber,style=authoryear-comp,uniquename=false,uniquelist=false]{biblatex}
\addbibresource{references.bib}


\usepackage{csquotes}
\usepackage{booktabs}
\usepackage{array}
\usepackage{multirow}
\usepackage{wrapfig}
\usepackage{float}
\usepackage{colortbl}
\usepackage{pdflscape}
\usepackage{threeparttable}
\usepackage{threeparttablex}
\usepackage[normalem]{ulem}
\usepackage{makecell}
\usepackage{xcolor}
\usepackage{multicol}
\usepackage{fancyhdr}
\usepackage{titlesec}
\usepackage{etoolbox}
\usepackage{tabularx}
\usepackage{afterpage}

% Define WISTA colors (brighter blue)
\definecolor{wistaBlue}{RGB}{30,110,180}
\definecolor{wistaLightBlue}{RGB}{51,153,204}
\definecolor{wistaDarkGray}{RGB}{64,64,64}
\definecolor{wistaLightGray}{RGB}{128,128,128}

% Custom chapter formatting - number on top with column-width rule  
\titleformat{\chapter}[display]
  {\normalfont\huge\bfseries\color{wistaDarkGray}}
  {\color{wistaBlue}\large\chaptertitlename\ \thechapter\vspace{0.2cm}\\\noindent\rule{\columnwidth}{1.5pt}}
  {0.5cm}
  {\Huge\color{wistaDarkGray}}
  [\vspace{0.2cm}\noindent{\color{wistaBlue}\rule{\columnwidth}{1pt}}\vspace{0.5cm}]
\titlespacing*{\chapter}{0pt}{-20pt}{20pt}

% Custom section formatting - number on top with column-width underline
\titleformat{\section}[display]
  {\normalfont\large\bfseries}
  {\color{wistaBlue}\large\bfseries\thesection\vspace{0.1cm}\\\noindent\rule{\columnwidth}{1pt}}
  {0.3cm}
  {\color{wistaDarkGray}\large\bfseries}
  [\vspace{0.1cm}\noindent{\color{wistaBlue}\rule{\columnwidth}{0.5pt}}\vspace{0.3cm}]
\titlespacing*{\section}{0pt}{1.5ex plus 1ex minus .2ex}{1ex plus .2ex}

% Smaller formatting for subsections (within columns)
\titleformat{\subsection}
  {\color{wistaDarkGray}\normalsize\bfseries}
  {\color{wistaBlue}\normalsize\bfseries\thesubsection}
  {0.5em}
  {\color{wistaDarkGray}\normalsize\bfseries}
  [\vspace{0.05cm}\noindent{\color{wistaBlue}\rule{0.5\columnwidth}{0.4pt}}\vspace{0.1cm}]
  
% Even smaller for subsubsections  
\titleformat{\subsubsection}
  {\color{wistaDarkGray}\small\bfseries}
  {\color{wistaBlue}\small\bfseries\thesubsubsection}
  {0.3em}
  {\color{wistaDarkGray}\small\bfseries}

% Header and footer setup
\pagestyle{fancy}
\fancyhf{} % Clear all header and footer fields
\fancyhead[LE]{\small\textit{Max Mustermann, Dr. Anna Schmidt}}
\fancyhead[RO]{\small\textit{Effizienzanalyse statistischer Verfahren}}
\fancyfoot[LE]{\small\thepage}
\fancyfoot[RE]{\small Statistisches Bundesamt | WISTA | {\color{wistaBlue}6} | 2025}
\fancyfoot[LO]{\small Statistisches Bundesamt | WISTA | {\color{wistaBlue}6} | 2025}
\fancyfoot[RO]{\small\thepage}
\renewcommand{\headrulewidth}{0pt}
\renewcommand{\footrulewidth}{0pt}

% Make all figures and tables span both columns (full width)
% Redefine figure environment to use figure*
\let\origfigure\figure
\let\endorigfigure\endfigure
\renewenvironment{figure}[1][htbp]{%
  \origfigure*[#1]%
}{%
  \endorigfigure*%
}

% Redefine table environment to use table*
\let\origtable\table
\let\endorigtable\endtable
\renewenvironment{table}[1][htbp]{%
  \origtable*[#1]%
}{%
  \endorigtable*%
}
\makeatletter
\@ifpackageloaded{bookmark}{}{\usepackage{bookmark}}
\makeatother
\makeatletter
\@ifpackageloaded{caption}{}{\usepackage{caption}}
\AtBeginDocument{%
\ifdefined\contentsname
  \renewcommand*\contentsname{Inhaltsverzeichnis}
\else
  \newcommand\contentsname{Inhaltsverzeichnis}
\fi
\ifdefined\listfigurename
  \renewcommand*\listfigurename{Abbildungsverzeichnis}
\else
  \newcommand\listfigurename{Abbildungsverzeichnis}
\fi
\ifdefined\listtablename
  \renewcommand*\listtablename{Tabellenverzeichnis}
\else
  \newcommand\listtablename{Tabellenverzeichnis}
\fi
\ifdefined\figurename
  \renewcommand*\figurename{Abbildung}
\else
  \newcommand\figurename{Abbildung}
\fi
\ifdefined\tablename
  \renewcommand*\tablename{Tabelle}
\else
  \newcommand\tablename{Tabelle}
\fi
}
\@ifpackageloaded{float}{}{\usepackage{float}}
\floatstyle{ruled}
\@ifundefined{c@chapter}{\newfloat{codelisting}{h}{lop}}{\newfloat{codelisting}{h}{lop}[chapter]}
\floatname{codelisting}{Listing}
\newcommand*\listoflistings{\listof{codelisting}{Listingverzeichnis}}
\makeatother
\makeatletter
\makeatother
\makeatletter
\@ifpackageloaded{caption}{}{\usepackage{caption}}
\@ifpackageloaded{subcaption}{}{\usepackage{subcaption}}
\makeatother
\usepackage{bookmark}
\IfFileExists{xurl.sty}{\usepackage{xurl}}{} % add URL line breaks if available
\urlstyle{same}
\hypersetup{
  pdftitle={Effizienzanalyse statistischer Verfahren in der amtlichen Statistik},
  pdfauthor={Max Mustermann; Dr.~Anna Schmidt},
  pdflang={de},
  hidelinks,
  pdfcreator={LaTeX via pandoc}}


\title{Effizienzanalyse statistischer Verfahren in der amtlichen
Statistik}
\usepackage{etoolbox}
\makeatletter
\providecommand{\subtitle}[1]{% add subtitle to \maketitle
  \apptocmd{\@title}{\par {\large #1 \par}}{}{}
}
\makeatother
\subtitle{Eine empirische Untersuchung am Beispiel der
Datenverarbeitung}
\author{Max Mustermann \and Dr.~Anna Schmidt}
\date{2025-11-22}
\begin{document}
\frontmatter
\maketitle

\renewcommand*\contentsname{Inhaltsverzeichnis}
{
\setcounter{tocdepth}{2}
\tableofcontents
}

\setstretch{1.1}
\mainmatter
\bookmarksetup{startatroot}

\chapter*{Einleitung}\label{einleitung}
\addcontentsline{toc}{chapter}{Einleitung}

\markboth{Einleitung}{Einleitung}

\section*{Zusammenfassung}\label{zusammenfassung}
\addcontentsline{toc}{section}{Zusammenfassung}

\markright{Zusammenfassung}

Die vorliegende Studie untersucht die Effizienz statistischer Verfahren
in der amtlichen Statistik am Beispiel moderner
Datenverarbeitungsmethoden. Dabei werden insbesondere die Auswirkungen
verschiedener algorithmischer Ansätze auf die Qualität und
Geschwindigkeit der statistischen Auswertungen analysiert. Die
Ergebnisse zeigen signifikante Verbesserungen bei der Verwendung
optimierter Verfahren.

\textbf{Schlagwörter:} Statistische Verfahren, Effizienzanalyse,
Datenverarbeitung, Amtliche Statistik, Algorithmusoptimierung

\section*{Abstract}\label{abstract}
\addcontentsline{toc}{section}{Abstract}

\markright{Abstract}

This study examines the efficiency of statistical procedures in official
statistics using modern data processing methods as an example. The
analysis focuses particularly on the effects of different algorithmic
approaches on the quality and speed of statistical evaluations. The
results show significant improvements when using optimized procedures.

\textbf{Keywords:} Statistical procedures, efficiency analysis, data
processing, official statistics, algorithm optimization

\section*{Motivation und Zielsetzung}\label{motivation-und-zielsetzung}
\addcontentsline{toc}{section}{Motivation und Zielsetzung}

\markright{Motivation und Zielsetzung}

Die amtliche Statistik steht vor der Herausforderung, bei stetig
wachsenden Datenmengen und sich ändernden Nutzeranforderungen qualitativ
hochwertige und zeitnahe statistische Informationen bereitzustellen
\autocite{destatis2023}. Dabei spielt die Effizienz der verwendeten
statistischen Verfahren eine zentrale Rolle für die Leistungsfähigkeit
des gesamten statistischen Systems.

In den vergangenen Jahren haben sich die Rahmenbedingungen für die
Produktion amtlicher Statistiken erheblich verändert. Die
Digitalisierung hat nicht nur zu einer Vervielfachung der verfügbaren
Datenquellen geführt, sondern auch neue Möglichkeiten für die
statistische Analyse eröffnet. Gleichzeitig sind die Erwartungen der
Nutzer an Aktualität und Detailliertheit der Statistiken gestiegen.

Diese Entwicklungen erfordern eine kontinuierliche Evaluation und
Optimierung der eingesetzten Verfahren. Ziel der vorliegenden Studie ist
es daher, die Effizienz verschiedener statistischer Verfahren
systematisch zu bewerten und Handlungsempfehlungen für die Praxis
abzuleiten.

\section*{Forschungsfragen}\label{forschungsfragen}
\addcontentsline{toc}{section}{Forschungsfragen}

\markright{Forschungsfragen}

Die Untersuchung fokussiert sich auf folgende zentrale Forschungsfragen:

\begin{enumerate}
\def\labelenumi{\arabic{enumi}.}
\item
  \textbf{Methodische Effizienz:} Wie unterscheiden sich verschiedene
  statistische Verfahren hinsichtlich ihrer Rechenzeit und
  Speicheranforderungen bei gleichbleibender Ergebnisqualität?
\item
  \textbf{Skalierbarkeit:} Welche Verfahren zeigen die beste Performance
  bei wachsenden Datenmengen, und wo liegen die kritischen Grenzen?
\item
  \textbf{Qualitätssicherung:} Inwiefern beeinflussen
  effizienzorientierte Optimierungen die statistische Qualität der
  Ergebnisse?
\end{enumerate}

\section*{Aufbau der Arbeit}\label{aufbau-der-arbeit}
\addcontentsline{toc}{section}{Aufbau der Arbeit}

\markright{Aufbau der Arbeit}

Die Arbeit gliedert sich wie folgt: Nach dieser Einleitung wird in
Kapitel~\ref{sec-btp} der Anwendungsfall Business-Tax-Panel vorgestellt.
Kapitel~\ref{sec-infrastruktur} beschreibt die Infrastruktur des
Forschungsdatenzentrums. In Kapitel~\ref{sec-methoden} werden
verschiedene Auswertungsmethoden vorgestellt.
Kapitel~\ref{sec-performance} präsentiert die empirischen
Performanzergebnisse. Die Arbeit schließt mit einem Ausblick in
Kapitel~\ref{sec-ausblick}.

Die Analyse basiert auf einer umfangreichen empirischen Studie, die
verschiedene statistische Verfahren anhand realer Datensätze aus der
amtlichen Statistik vergleicht. Dabei werden sowohl etablierte als auch
neuere Ansätze berücksichtigt, um ein vollständiges Bild der verfügbaren
Optionen zu zeichnen.

\bookmarksetup{startatroot}

\chapter{Einleitung}\label{sec-einleitung}

\section{Motivation und Zielsetzung}\label{motivation-und-zielsetzung-1}

Die amtliche Statistik steht vor der Herausforderung, bei stetig
wachsenden Datenmengen und sich ändernden Nutzeranforderungen qualitativ
hochwertige und zeitnahe statistische Informationen bereitzustellen
\autocite{destatis2023}. Dabei spielt die Effizienz der verwendeten
statistischen Verfahren eine zentrale Rolle für die Leistungsfähigkeit
des gesamten statistischen Systems.

In den vergangenen Jahren haben sich die Rahmenbedingungen für die
Produktion amtlicher Statistiken erheblich verändert. Die
Digitalisierung hat nicht nur zu einer Vervielfachung der verfügbaren
Datenquellen geführt, sondern auch neue Möglichkeiten für die
statistische Analyse eröffnet. Gleichzeitig sind die Erwartungen der
Nutzer an Aktualität und Detailliertheit der Statistiken gestiegen.

Diese Entwicklungen erfordern eine kontinuierliche Evaluation und
Optimierung der eingesetzten Verfahren. Ziel der vorliegenden Studie ist
es daher, die Effizienz verschiedener statistischer Verfahren
systematisch zu bewerten und Handlungsempfehlungen für die Praxis
abzuleiten.

\section{Forschungsfragen}\label{forschungsfragen-1}

Die Untersuchung fokussiert sich auf folgende zentrale Forschungsfragen:

\begin{enumerate}
\def\labelenumi{\arabic{enumi}.}
\item
  \textbf{Methodische Effizienz:} Wie unterscheiden sich verschiedene
  statistische Verfahren hinsichtlich ihrer Rechenzeit und
  Speicheranforderungen bei gleichbleibender Ergebnisqualität?
\item
  \textbf{Skalierbarkeit:} Welche Verfahren zeigen die beste Performance
  bei wachsenden Datenmengen, und wo liegen die kritischen Grenzen?
\item
  \textbf{Praktische Anwendbarkeit:} Welche Infrastrukturanforderungen
  stellen moderne Auswertungsmethoden an Forschungsdatenzentren?
\end{enumerate}

\section{Aufbau der Arbeit}\label{aufbau-der-arbeit-1}

Die Arbeit gliedert sich wie folgt: Nach dieser Einleitung wird in
Kapitel~\ref{sec-btp} der Anwendungsfall Business-Tax-Panel vorgestellt,
der als empirisches Beispiel für die Analysen dient.
Kapitel~\ref{sec-infrastruktur} beschreibt die Infrastruktur des
Forschungsdatenzentrums des Bundes (FDZ-Bund), in der die Auswertungen
durchgeführt werden.

In Kapitel~\ref{sec-methoden} werden verschiedene effiziente
Auswertungsmethoden vorgestellt und deren theoretische Grundlagen
erläutert. Kapitel~\ref{sec-performance} präsentiert die empirischen
Performanzergebnisse der verschiedenen Ansätze. Die Arbeit schließt mit
einem Ausblick auf zukünftige Entwicklungen in
Kapitel~\ref{sec-ausblick}.

Die Analyse basiert auf einer umfangreichen empirischen Studie, die
verschiedene statistische Verfahren anhand realer Datensätze aus der
amtlichen Statistik vergleicht. Dabei werden sowohl etablierte als auch
neuere Ansätze berücksichtigt, um ein vollständiges Bild der verfügbaren
Optionen zu zeichnen.

\bookmarksetup{startatroot}

\chapter{Anwendungsfall Business-Tax-Panel}\label{sec-btp}

\section{Überblick zum
Business-Tax-Panel}\label{uxfcberblick-zum-business-tax-panel}

Das Business-Tax-Panel (BTP) ist ein zentraler Forschungsdatensatz des
Statistischen Bundesamtes, der umfassende Informationen zur
Unternehmensbesteuerung in Deutschland bereitstellt. Das Panel
kombiniert Daten aus verschiedenen Steuerstatistiken miteinander und
ermöglicht somit eine umfassende Analyse der Besteuerung von Unternehmen
über Zeit.

Im BTP werden folgende Datensätze zusammengeführt:

\begin{itemize}
\tightlist
\item
  \textbf{Gewerbesteuerstatistik} (EVAS-Nr. 73511)
\item
  \textbf{Körperschaftsteuerstatistik} (EVAS-Nr. 73211)
\item
  \textbf{Umsatzsteuerstatistik - Voranmeldungen} (EVAS-Nr. 73311)
\item
  \textbf{Statistik über Personengesellschaften und Gemeinschaften}
  (EVAS-Nr. 73121)
\item
  \textbf{Umsatzsteuerstatistik - Veranlagungen} (EVAS-Nr. 73321)
\item
  \textbf{Einnahmenüberschussrechnung} aus der Lohn- und
  Einkommensteuerstatistik (EVAS-Nr. 73111)
\item
  Merkmale aus dem \textbf{Statistischen Unternehmensregister} (EVAS-Nr.
  52111)
\end{itemize}

\section{Datenstruktur und Umfang}\label{datenstruktur-und-umfang}

\subsection{Dimensionen des
Datensatzes}\label{dimensionen-des-datensatzes}

Das BTP umfasst einen außergewöhnlich umfangreichen Datensatz mit
beeindruckenden Dimensionen:

\begin{table}

\caption{\label{tbl-btp-overview}Übersicht über die Dimensionen des
Business-Tax-Panel}

\centering{

\begin{tabular}{ll}
\toprule
Merkmal & Wert\\
\midrule
Beobachtungszeitraum & 2013-2019 (jährliche Erweiterung)\\
Anzahl Einheiten (2013-2019) & >66 Millionen\\
Anzahl Merkmale & >2.700\\
Erhebungsart & Vollerhebung\\
Datenquelle & Sekundärerhebung (Finanzverwaltung)\\
\bottomrule
\end{tabular}

}

\end{table}%

\subsection{Hauptmerkmale des
Datensatzes}\label{hauptmerkmale-des-datensatzes}

Aus den verschiedenen Steuervordrucken und darin verwendeten Kennzahlen
werden über 2.700 Merkmale erfasst, die verschiedene Aspekte der
Unternehmensbesteuerung abbilden:

\begin{itemize}
\tightlist
\item
  \textbf{Steuerliche Kennzahlen:} Bemessungsgrundlagen, Steuerbeträge,
  Vorsteuern aus verschiedenen Steuerarten
\item
  \textbf{Wirtschaftliche Kennzahlen:} Umsätze, Gewinne,
  Betriebsausgaben
\item
  \textbf{Strukturmerkmale:} Rechtsform, Wirtschaftszweig, Größenklassen
\item
  \textbf{Geografische Informationen:} Bundesland, Gemeindekennziffer
\item
  \textbf{Panelstruktur:} Längsschnittinformationen ermöglichen
  Zeitreihenanalysen
\end{itemize}

\subsection{Datenerhebung und
Rechtsgrundlage}\label{datenerhebung-und-rechtsgrundlage}

Die im BTP enthaltenen Steuerstatistiken sind Bundesstatistiken, die als
\textbf{Vollerhebung} durchgeführt werden. Es handelt sich um
\textbf{Sekundärerhebungen}: Die Erhebungsmerkmale werden aus den
Voranmeldungs- und Veranlagungsbescheiden von der Finanzverwaltung
entnommen und über deren Rechenzentren an die Statistischen Ämter der
Länder übermittelt.

Rechtsgrundlage ist das
\href{https://www.gesetze-im-internet.de/ststatg_1995/}{Gesetz über
Steuerstatistiken (StStatG)} vom 11. Oktober 1995.

\section{Forschungspotenzial und
Nutzung}\label{forschungspotenzial-und-nutzung}

\subsection{Typische Forschungsfragen}\label{typische-forschungsfragen}

Das BTP wird für eine Vielzahl von Forschungsfragen genutzt:

\begin{enumerate}
\def\labelenumi{\arabic{enumi}.}
\tightlist
\item
  \textbf{Steuerliche Analysen:} Effekte von Steuerreformen auf
  Unternehmensentscheidungen
\item
  \textbf{Unternehmensdynamik:} Gründungen, Wachstum und Marktaustritte
\item
  \textbf{Wirtschaftsstruktur:} Branchenentwicklung und regionale
  Unterschiede
\item
  \textbf{Krisenanalysen:} Auswirkungen von wirtschaftlichen Schocks auf
  Unternehmen
\end{enumerate}

\subsection{Herausforderungen bei der
Datenauswertung}\label{herausforderungen-bei-der-datenauswertung}

Die Größe und Komplexität des BTP stellt besondere Anforderungen an die
Auswertungsinfrastruktur:

\begin{itemize}
\tightlist
\item
  \textbf{Rechenleistung:} Analysen über den gesamten Datensatz
  erfordern erhebliche Rechenkapazitäten
\item
  \textbf{Speicherbedarf:} Das Vorhalten des Datensatzes im
  Arbeitsspeicher ist oft nicht möglich
\item
  \textbf{Komplexe Verknüpfungen:} Die Zusammenführung verschiedener
  Datenquellen erfordert effiziente Join-Operationen
\item
  \textbf{Datenschutz:} Strenge Anforderungen an den Umgang mit
  sensiblen Unternehmensdaten
\end{itemize}

\section{Zugangsformen im FDZ}\label{zugangsformen-im-fdz}

Das BTP ist über verschiedene Zugangsformen verfügbar:

\begin{itemize}
\tightlist
\item
  \textbf{Gastwissenschaftsarbeitsplatz (GWAP):} On-Site-Zugang mit
  vollständigen Daten, auswertbar mit SAS oder Stata. Bedingt durch die
  Speicherplatzgröße kann die Nutzung einer Stichprobe erforderlich
  sein.
\item
  \textbf{Kontrollierte Datenfernverarbeitung (KDFV):} Remote-Zugang mit
  Ergebnisfreigabe über KDFV
\item
  Die Gemeindekennziffer für Bayern steht am GWAP lediglich
  pseudonymisiert zur Verfügung
\end{itemize}

\section{BTP als Benchmark-Datensatz}\label{btp-als-benchmark-datensatz}

Aufgrund seiner Größe (\textgreater66 Millionen Einheiten), Komplexität
(\textgreater2.700 Variablen) und realen Datenstruktur eignet sich das
BTP hervorragend als Benchmark für die Evaluation verschiedener
Auswertungsmethoden. Die Herausforderungen des BTP sind repräsentativ
für viele Forschungsprojekte im FDZ:

\begin{itemize}
\tightlist
\item
  \textbf{Datengröße:} Der Datensatz übersteigt oft die verfügbare
  RAM-Kapazität
\item
  \textbf{Komplexe Joins:} Verknüpfung verschiedener Steuerstatistiken
  erforderlich\\
\item
  \textbf{Aggregationen:} Über verschiedene Dimensionen (Zeit, Region,
  Wirtschaftszweig)
\item
  \textbf{Filteroperationen:} Selektion relevanter Unternehmensgruppen
\item
  \textbf{Zeitreihenanalysen:} Panel-Struktur über mehrere Jahre
\end{itemize}

Diese realitätsnahen Szenarien ermöglichen es, die Praxistauglichkeit
verschiedener Auswertungsmethoden unter den tatsächlichen Bedingungen
eines Forschungsdatenzentrums zu bewerten.

\bookmarksetup{startatroot}

\chapter{Infrastruktur des FDZ-Bund}\label{sec-infrastruktur}

\section{Die Forschungsdatenzentren der amtlichen
Statistik}\label{die-forschungsdatenzentren-der-amtlichen-statistik}

Die Forschungsdatenzentren (FDZ) der amtlichen Statistik bestehen aus
zwei organisatorisch voneinander getrennten Einrichtungen: dem
\textbf{Forschungsdatenzentrum des Statistischen Bundesamts} und dem
\textbf{Forschungsdatenzentrum der Statistischen Ämter der Länder}.
Beide FDZ arbeiten eng zusammen und bieten ein abgestimmtes Daten- und
Dienstleistungsangebot für die wissenschaftliche Nutzung von Mikrodaten
der amtlichen Statistiken.

\subsection{Entstehung und
Entwicklung}\label{entstehung-und-entwicklung}

Die Gründung der FDZ geht auf ein Gutachten der „Kommission zur
Verbesserung der informationellen Infrastruktur zwischen Wissenschaft
und Statistik'' (KVI) von 1999 zurück. Mit finanzieller Förderung des
Bundesministeriums für Bildung und Forschung (BMBF) wurde im Herbst 2001
zunächst das FDZ des Statistischen Bundesamts gegründet. Das FDZ der
Statistischen Ämter der Länder folgte im April 2002.

\subsection{Aufgaben und Ziele}\label{aufgaben-und-ziele}

Die FDZ verfolgen das Ziel, den Zugang zu Mikrodaten der amtlichen
Statistik für die Wissenschaft nicht nur zu ermöglichen, sondern
fortlaufend zu verbessern und an die sich ändernden Bedarfe anzupassen:

\begin{itemize}
\tightlist
\item
  \textbf{Datenzugang:} Bereitstellung hochwertiger Mikrodaten für die
  wissenschaftliche Forschung
\item
  \textbf{Datenschutz:} Gewährleistung des Schutzes personenbezogener
  und unternehmensbezogener Daten nach § 16 BStatG
\item
  \textbf{Beratung:} Unterstützung der Forschenden bei der Datennutzung
  und methodischen Fragen
\item
  \textbf{Infrastruktur:} Bereitstellung einer sicheren und
  leistungsfähigen Auswertungsumgebung
\item
  \textbf{Regionalisierung:} Flächendeckende Standorte in ganz
  Deutschland für kurze Anfahrtswege
\end{itemize}

\section{Technische Infrastruktur}\label{technische-infrastruktur}

\subsection{Hardware-Ausstattung}\label{hardware-ausstattung}

Die technische Infrastruktur des FDZ-Bund umfasst verschiedene
Systemkomponenten:

\begin{table}

\caption{\label{tbl-fdz-infrastructure}Technische Ausstattung des
FDZ-Bund (beispielhafte Konfiguration)}

\centering{

\begin{tabular}{llc}
\toprule
Komponente & Spezifikation & Anzahl\\
\midrule
Rechenknoten (Standard) & 2x Intel Xeon Gold 6248R (48 Kerne), 384 GB RAM & 12\\
Rechenknoten (High-Memory) & 2x Intel Xeon Platinum 8280 (56 Kerne), 1,5 TB RAM & 4\\
Speichersystem & NAS-System mit 500 TB nutzbarem Speicher & 1\\
Netzwerk & 10 Gbit/s Ethernet, InfiniBand für parallele Operationen & -\\
Backup-System & Tägliches Backup, räumlich getrennte Archivierung & -\\
\bottomrule
\end{tabular}

}

\end{table}%

\subsection{Software-Umgebung}\label{software-umgebung}

Den Forschenden stehen verschiedene Software-Tools zur Verfügung:

\begin{itemize}
\tightlist
\item
  \textbf{Statistische Software:} R, Python, Stata, SAS, SPSS
\item
  \textbf{Datenbanksysteme:} PostgreSQL, DuckDB
\item
  \textbf{Entwicklungsumgebungen:} RStudio Server, JupyterLab
\item
  \textbf{Versionskontrolle:} Git (lokal, ohne externe Verbindung)
\end{itemize}

\section{Infrastrukturkonzept}\label{infrastrukturkonzept}

\subsection{Fachlich zentralisierte
Datenhaltung}\label{fachlich-zentralisierte-datenhaltung}

In Deutschland wird der überwiegende Teil der amtlichen
Statistikproduktion dezentral in den Statistischen Ämtern der Länder
durchgeführt (über 90\% aller amtlichen Statistiken). Da sich
wissenschaftliche Analysen in der Regel jedoch auf mehrere Bundesländer
oder das gesamte Bundesgebiet beziehen, haben die FDZ eine
\textbf{fachlich zentralisierte Datenhaltung} eingerichtet. Diese
ermöglicht es, die Daten des gesamten Bundesgebiets länderübergreifend
an jedem Standort bereitzustellen.

\subsection{Regionalisierte
Infrastruktur}\label{regionalisierte-infrastruktur}

Die FDZ verfolgen das Ziel, immer in der Nähe der Wissenschaft zu sein.
Dies wird durch über ganz Deutschland verteilte Standorte erreicht.
Standorte befinden sich u.a. in:

\begin{itemize}
\tightlist
\item
  Berlin, Bonn, Wiesbaden (Bund)
\item
  Bad Ems, Bremen, Dresden, Düsseldorf, Erfurt, Frankfurt, Hamburg,
  Hannover, Kiel, München, Stuttgart u.v.m. (Länder)
\end{itemize}

\section{Zugangsformen und
Nutzungsmodelle}\label{zugangsformen-und-nutzungsmodelle}

Die FDZ bieten verschiedene Zugangswege an, über die unterschiedlich
anonymisierte Datenprodukte bereitgestellt werden:

\subsection{Off-Site-Zugang}\label{off-site-zugang}

\begin{itemize}
\tightlist
\item
  \textbf{Public Use Files (PUF):} Stark anonymisierte Daten für freie
  Nutzung
\item
  \textbf{Scientific Use Files (SUF):} Faktisch anonymisierte Daten für
  wissenschaftliche Zwecke
\item
  \textbf{Remote Scientific Use Files (Remote-SUF):} Neuer Zugangsweg
  mit höherer Detailtiefe
\end{itemize}

\subsection{On-Site-Zugang}\label{on-site-zugang}

\begin{itemize}
\tightlist
\item
  \textbf{Gastwissenschaftsarbeitsplatz (GWAP):} Zugang zu formal
  anonymisierten Daten vor Ort

  \begin{itemize}
  \tightlist
  \item
    Vorteile: Höchste Datendetailtiefe, direkte Beratung
  \item
    Anforderungen: Terminvereinbarung, Zutrittskontrolle,
    Output-Kontrolle
  \end{itemize}
\item
  \textbf{Kontrollierte Datenfernverarbeitung (KDFV):} Remote-Ausführung
  von Analyseskripten

  \begin{itemize}
  \tightlist
  \item
    Vorteile: Keine Anreise, zeitlich flexibel
  \item
    Einschränkungen: Keine interaktive Arbeit, Ergebnisfreigabe
    erforderlich
  \end{itemize}
\end{itemize}

\section{Datenschutz und Sicherheit}\label{datenschutz-und-sicherheit}

\subsection{Mehrstufiges
Sicherheitskonzept}\label{mehrstufiges-sicherheitskonzept}

Das FDZ-Bund implementiert ein umfassendes Sicherheitskonzept:

\begin{enumerate}
\def\labelenumi{\arabic{enumi}.}
\tightlist
\item
  \textbf{Physische Sicherheit:} Zutrittskontrolle, gesicherte
  Räumlichkeiten
\item
  \textbf{Netzwerksicherheit:} Isolierte Netzwerke, Firewalls, Intrusion
  Detection
\item
  \textbf{Zugriffsrechte:} Rollenbasierte Zugriffskontrolle,
  Authentifizierung
\item
  \textbf{Output-Kontrolle:} Prüfung aller Ergebnisse vor Freigabe
\item
  \textbf{Audit-Logging:} Protokollierung aller Zugriffe und Operationen
\end{enumerate}

\subsection{Output-Kontrolle}\label{output-kontrolle}

Bevor Forschungsergebnisse das FDZ verlassen dürfen, durchlaufen sie
eine mehrstufige Kontrolle:

\begin{itemize}
\tightlist
\item
  \textbf{Automatische Prüfung:} Screening auf kritische Werte (z.B.
  Fallzahlen \textless{} 5)
\item
  \textbf{Manuelle Prüfung:} Sichtung durch geschulte Mitarbeiter
\item
  \textbf{Dokumentation:} Nachvollziehbare Begründung der
  Freigabeentscheidung
\end{itemize}

\section{Performanz-Optimierung in der
FDZ-Umgebung}\label{performanz-optimierung-in-der-fdz-umgebung}

\subsection{Herausforderungen}\label{herausforderungen}

Die Arbeit in der geschützten FDZ-Umgebung stellt besondere
Anforderungen:

\begin{itemize}
\tightlist
\item
  \textbf{Begrenzte Ressourcen:} Mehrere Nutzer teilen sich die
  verfügbare Hardware
\item
  \textbf{Keine Cloud-Dienste:} Externe Rechenressourcen können nicht
  genutzt werden
\item
  \textbf{Datengröße:} Große Datensätze müssen effizient verarbeitet
  werden
\item
  \textbf{Reproduzierbarkeit:} Analysen müssen dokumentiert und
  nachvollziehbar sein
\end{itemize}

\subsection{Optimierungsstrategien}\label{optimierungsstrategien}

Verschiedene Strategien helfen, die Performanz zu optimieren:

\begin{figure}

\centering{

\pandocbounded{\includegraphics[keepaspectratio]{sections/03-infrastruktur_files/figure-pdf/fig-optimization-strategies-1.pdf}}

}

\caption{\label{fig-optimization-strategies}Strategien zur
Performanz-Optimierung im FDZ-Bund}

\end{figure}%

\section{Akkreditierung und
Qualitätssicherung}\label{akkreditierung-und-qualituxe4tssicherung}

Die Forschungsdatenzentren der Statistischen Ämter des Bundes und der
Länder sind durch den \textbf{Rat für Sozial- und WirtschaftsDaten
(RatSWD) akkreditiert}. Der RatSWD akkreditiert seit 2008
Forschungsdatenzentren, die den Zugang zu qualitativ hochwertigen
sensiblen Daten für die Wissenschaft ermöglichen und dabei die
erforderlichen Kriterien erfüllen.

\section{Best Practices für
FDZ-Nutzende}\label{best-practices-fuxfcr-fdz-nutzende}

Aus der Erfahrung mit den FDZ lassen sich folgende Empfehlungen für
effizientes Arbeiten ableiten:

\begin{enumerate}
\def\labelenumi{\arabic{enumi}.}
\tightlist
\item
  \textbf{Datenreduktion:} Frühzeitige Filterung auf relevante
  Beobachtungen und Variablen
\item
  \textbf{Inkrementelle Entwicklung:} Erst mit Stichproben arbeiten,
  dann auf Vollsatz skalieren
\item
  \textbf{Effiziente Datenformate:} Nutzung moderner Formate
  (Arrow/Parquet) statt CSV
\item
  \textbf{Dokumentation:} Sorgfältige Dokumentation des Analyseprozesses
  für Output-Kontrolle
\item
  \textbf{Ressourcenplanung:} Bewusste Planung rechenintensiver
  Operationen bei geteilter Infrastruktur
\item
  \textbf{Beratung nutzen:} Frühzeitige Kontaktaufnahme mit
  FDZ-Mitarbeitern bei methodischen Fragen
\end{enumerate}

\bookmarksetup{startatroot}

\chapter{Effiziente Auswertungsmethoden}\label{sec-methoden}

\section{Motivation}\label{motivation}

Die Analyse großer Datensätze wie des BTP (\textgreater66 Millionen
Einheiten, \textgreater2.700 Variablen) in der FDZ-Umgebung stellt
besondere Anforderungen an die verwendeten Werkzeuge. Traditionelle
Ansätze (CSV-Dateien, Base R) stoßen bei dieser Datengröße an ihre
Grenzen. Moderne spaltenorientierte Technologien bieten hier deutliche
Vorteile.

\section{Apache Arrow: Spaltenorientiertes
In-Memory-Format}\label{apache-arrow-spaltenorientiertes-in-memory-format}

\subsection{Grundprinzip}\label{grundprinzip}

Apache Arrow ist ein sprach- und plattformübergreifendes
In-Memory-Format für spaltenorientierte Datenverarbeitung. Im Gegensatz
zu zeilenorientierten Formaten werden Daten spaltenweise organisiert,
was erhebliche Vorteile bietet:

\begin{itemize}
\tightlist
\item
  \textbf{Effiziente Kompression:} Ähnliche Werte in einer Spalte lassen
  sich besser komprimieren
\item
  \textbf{Schnellere Aggregationen:} Nur benötigte Spalten müssen
  gelesen werden (Columnar Scanning)
\item
  \textbf{Zero-Copy:} Daten können ohne Kopieren zwischen Prozessen und
  Sprachen ausgetauscht werden
\item
  \textbf{SIMD-Optimierung:} Moderne CPU-Instruktionen können effizient
  genutzt werden
\end{itemize}

\subsection{Anwendung in R}\label{anwendung-in-r}

Das \texttt{arrow}-Paket in R bietet vollständige Integration:

\begin{Shaded}
\begin{Highlighting}[]
\FunctionTok{library}\NormalTok{(arrow)}

\CommentTok{\# Lesen eines Parquet{-}Datasets mit Filterung}
\NormalTok{dataset }\OtherTok{\textless{}{-}} \FunctionTok{open\_dataset}\NormalTok{(}\StringTok{"btp\_data.parquet"}\NormalTok{)}

\NormalTok{result }\OtherTok{\textless{}{-}}\NormalTok{ dataset }\SpecialCharTok{|\textgreater{}}
  \FunctionTok{filter}\NormalTok{(year }\SpecialCharTok{\textgreater{}=} \DecValTok{2015}\NormalTok{, sector }\SpecialCharTok{==} \StringTok{"Manufacturing"}\NormalTok{) }\SpecialCharTok{|\textgreater{}}
  \FunctionTok{select}\NormalTok{(company\_id, year, revenue, employees) }\SpecialCharTok{|\textgreater{}}
  \FunctionTok{group\_by}\NormalTok{(year) }\SpecialCharTok{|\textgreater{}}
  \FunctionTok{summarise}\NormalTok{(}
    \AttributeTok{total\_revenue =} \FunctionTok{sum}\NormalTok{(revenue),}
    \AttributeTok{mean\_employees =} \FunctionTok{mean}\NormalTok{(employees)}
\NormalTok{  ) }\SpecialCharTok{|\textgreater{}}
  \FunctionTok{collect}\NormalTok{()  }\CommentTok{\# Erst hier wird die Berechnung ausgeführt}
\end{Highlighting}
\end{Shaded}

\subsection{Vorteile für
FDZ-Analysen}\label{vorteile-fuxfcr-fdz-analysen}

\begin{itemize}
\tightlist
\item
  \textbf{Out-of-Core Processing:} Datensätze müssen nicht vollständig
  in den RAM passen
\item
  \textbf{Lazy Evaluation:} Operationen werden optimiert, bevor sie
  ausgeführt werden
\item
  \textbf{Predicate Pushdown:} Filter werden direkt beim Lesen
  angewendet
\item
  \textbf{Partitionierung:} Datensätze können nach Kriterien (z.B. Jahr)
  partitioniert werden
\end{itemize}

\subsection{Limitationen}\label{limitationen}

\begin{itemize}
\tightlist
\item
  \textbf{Lernkurve:} Neue Syntax für Nutzer mit Base-R-Erfahrung
\item
  \textbf{Noch in Entwicklung:} Nicht alle R-Funktionen werden
  unterstützt
\item
  \textbf{Setup-Aufwand:} Daten müssen einmalig ins Parquet-Format
  konvertiert werden
\end{itemize}

\section{Parquet: Persistentes spaltenorientiertes
Format}\label{parquet-persistentes-spaltenorientiertes-format}

\subsection{Eigenschaften}\label{eigenschaften}

Parquet ist ein spaltenorientiertes Speicherformat, das optimal mit
Arrow zusammenarbeitet:

\begin{itemize}
\tightlist
\item
  \textbf{Kompression:} Typischerweise 5-10x kleinere Dateien als CSV
\item
  \textbf{Metadaten:} Enthält Schema-Informationen und Statistiken
\item
  \textbf{Partitionierung:} Ermöglicht selektives Lesen nach
  Partitionskeys
\item
  \textbf{Kompatibilität:} Von vielen Tools unterstützt (R, Python,
  Spark, DuckDB)
\end{itemize}

\subsection{Vergleich zu Alternativen}\label{vergleich-zu-alternativen}

\begin{table}

\caption{\label{tbl-format-comparison}Vergleich verschiedener
Datenformate für die BTP-Analyse}

\centering{

\begin{tabular}{lrllcl}
\toprule
Format & Größe & Ladezeit & Kompression & Partitionierung & Kompatibilität\\
\midrule
CSV & 100\% & Sehr langsam & Keine & Nein & Universal\\
RDS (R) & 45\% & Mittel & Gzip & Nein & Nur R\\
Feather & 95\% & Sehr schnell & LZ4/ZSTD & Nein & Arrow-Ökosystem\\
Parquet & 15\% & Schnell & Snappy/GZIP & Möglich & Sehr hoch\\
Parquet (partitioniert) & 15\% & Sehr schnell (selektiv) & Snappy/GZIP & Ja & Sehr hoch\\
\bottomrule
\end{tabular}

}

\end{table}%

\subsection{Empfehlung für FDZ}\label{empfehlung-fuxfcr-fdz}

\textbf{Parquet ist das empfohlene Format für große Datensätze im FDZ}
wegen:

\begin{enumerate}
\def\labelenumi{\arabic{enumi}.}
\tightlist
\item
  Drastischer Reduktion der Dateigröße (wichtig bei begrenztem
  Speicherplatz)
\item
  Deutlich schnelleren Ladezeiten
\item
  Möglichkeit der selektiven Filterung ohne vollständiges Laden
\item
  Kompatibilität mit verschiedenen Analysewerkzeugen
\end{enumerate}

\section{DuckDB: Analytical Database
System}\label{duckdb-analytical-database-system}

\subsection{Konzept}\label{konzept}

DuckDB ist ein einbettbares analytisches Datenbanksystem, speziell für
OLAP-Workloads (Online Analytical Processing) entwickelt:

\begin{itemize}
\tightlist
\item
  \textbf{In-Process:} Läuft im gleichen Prozess wie R (keine separate
  Datenbank-Server-Installation nötig)
\item
  \textbf{Columnar Vectorized:} Spaltenorientierte Verarbeitung in
  SIMD-optimierten Batches
\item
  \textbf{SQL-Interface:} Bekannte SQL-Syntax für komplexe Abfragen
\item
  \textbf{Arrow-Integration:} Zero-Copy-Integration mit Arrow-Daten
\end{itemize}

\subsection{Warum DuckDB?}\label{warum-duckdb}

DuckDB löst spezifische Probleme bei der Analyse großer Datensätze:

\begin{enumerate}
\def\labelenumi{\arabic{enumi}.}
\tightlist
\item
  \textbf{Komplexe Aggregationen:} Window Functions, CTEs, komplexe
  Joins
\item
  \textbf{SQL-Komfort:} Mächtige Abfragesprache statt R-Code für
  Datentransformationen
\item
  \textbf{Out-of-Core:} Kann größere Daten als RAM verarbeiten
\item
  \textbf{Geschwindigkeit:} Hochoptimierte Query Engine
\end{enumerate}

\subsection{Anwendungsbeispiel}\label{anwendungsbeispiel}

\begin{Shaded}
\begin{Highlighting}[]
\FunctionTok{library}\NormalTok{(duckdb)}
\FunctionTok{library}\NormalTok{(arrow)}

\CommentTok{\# Verbindung erstellen}
\NormalTok{con }\OtherTok{\textless{}{-}} \FunctionTok{dbConnect}\NormalTok{(duckdb}\SpecialCharTok{::}\FunctionTok{duckdb}\NormalTok{())}

\CommentTok{\# Arrow{-}Dataset registrieren (ohne Laden)}
\FunctionTok{duckdb\_register\_arrow}\NormalTok{(con, }\StringTok{"btp"}\NormalTok{, arrow\_dataset)}

\CommentTok{\# Komplexe SQL{-}Abfrage}
\NormalTok{result }\OtherTok{\textless{}{-}} \FunctionTok{dbGetQuery}\NormalTok{(con, }\StringTok{"}
\StringTok{  SELECT }
\StringTok{    year,}
\StringTok{    sector,}
\StringTok{    COUNT(*) as n\_companies,}
\StringTok{    MEDIAN(revenue) as median\_revenue,}
\StringTok{    PERCENTILE\_CONT(0.9) WITHIN GROUP (ORDER BY revenue) as p90\_revenue,}
\StringTok{    LAG(MEDIAN(revenue)) OVER (}
\StringTok{      PARTITION BY sector }
\StringTok{      ORDER BY year}
\StringTok{    ) as prev\_year\_median}
\StringTok{  FROM btp}
\StringTok{  WHERE year BETWEEN 2015 AND 2019}
\StringTok{    AND revenue \textgreater{} 0}
\StringTok{  GROUP BY year, sector}
\StringTok{  ORDER BY year, sector}
\StringTok{"}\NormalTok{)}

\FunctionTok{dbDisconnect}\NormalTok{(con, }\AttributeTok{shutdown =} \ConstantTok{TRUE}\NormalTok{)}
\end{Highlighting}
\end{Shaded}

\subsection{Vorteile}\label{vorteile}

\begin{itemize}
\tightlist
\item
  \textbf{Kein Setup:} Keine separate Datenbank-Installation
  erforderlich
\item
  \textbf{Hohe Performance:} Oft schneller als spezialisierte
  Data-Frame-Bibliotheken
\item
  \textbf{SQL-Funktionen:} Umfangreiche Funktionen (Window Functions,
  Array-Operationen)
\item
  \textbf{Ressourceneffizient:} Intelligentes Memory Management
\item
  \textbf{R-Integration:} Direktes Arbeiten mit Arrow-Datasets oder
  R-Data-Frames
\end{itemize}

\subsection{Limitationen}\label{limitationen-1}

\begin{itemize}
\tightlist
\item
  \textbf{SQL-Kenntnisse:} Erfordert SQL-Grundkenntnisse
\item
  \textbf{Debugging:} SQL-Fehler können schwerer zu debuggen sein als
  R-Code
\item
  \textbf{Noch jung:} Weniger etabliert als traditionelle Datenbanken
  (aber sehr aktiv entwickelt)
\end{itemize}

\section{Vergleich mit Alternativen}\label{vergleich-mit-alternativen}

\subsection{Base R}\label{base-r}

\textbf{Vorteile:} Einfach, weit verbreitet, keine zusätzlichen
Abhängigkeiten \textbf{Nachteile:} Langsam bei großen Daten, hoher
Speicherverbrauch, alle Daten müssen in RAM passen

\textbf{Empfehlung:} Nur für kleine Datensätze (\textless1 GB) oder wenn
Kompatibilität wichtiger als Performance ist

\subsection{data.table}\label{data.table}

\textbf{Vorteile:} Sehr schnell für In-Memory-Operationen, ausgereift,
kompakte Syntax \textbf{Nachteile:} Alle Daten müssen in RAM passen,
steile Lernkurve, keine Out-of-Core-Unterstützung

\textbf{Empfehlung:} Gute Alternative wenn Datensatz in RAM passt und
keine Parquet-Integration benötigt wird

\subsection{dplyr + dbplyr}\label{dplyr-dbplyr}

\textbf{Vorteile:} Vertraute dplyr-Syntax, gut mit Datenbanken
integriert \textbf{Nachteile:} Langsamer als Arrow/DuckDB, benötigt
Backend-Datenbank, SQL-Translation manchmal limitiert

\textbf{Empfehlung:} Wenn bereits PostgreSQL- oder ähnliche
Datenbank-Infrastruktur vorhanden ist

\section{Unsere Empfehlung: Arrow + DuckDB +
Parquet}\label{unsere-empfehlung-arrow-duckdb-parquet}

Für die Analyse großer Datensätze im FDZ empfehlen wir die Kombination:

\begin{enumerate}
\def\labelenumi{\arabic{enumi}.}
\tightlist
\item
  \textbf{Datenspeicherung:} Parquet-Format, idealerweise partitioniert
  nach relevanten Dimensionen (z.B. Jahr)
\item
  \textbf{Datenmanipulation:} Arrow für einfache
  Filter/Select-Operationen
\item
  \textbf{Komplexe Analysen:} DuckDB für Aggregationen, Joins, Window
  Functions
\item
  \textbf{Ergebnisse:} R-Data-Frames für finale Visualisierung und
  Modellierung
\end{enumerate}

\subsection{Workflow-Beispiel}\label{workflow-beispiel}

\begin{Shaded}
\begin{Highlighting}[]
\FunctionTok{library}\NormalTok{(arrow)}
\FunctionTok{library}\NormalTok{(duckdb)}

\CommentTok{\# 1. Daten als Parquet{-}Dataset öffnen}
\NormalTok{btp }\OtherTok{\textless{}{-}} \FunctionTok{open\_dataset}\NormalTok{(}\StringTok{"data/btp.parquet"}\NormalTok{)}

\CommentTok{\# 2. Einfache Filterung mit Arrow}
\NormalTok{filtered }\OtherTok{\textless{}{-}}\NormalTok{ btp }\SpecialCharTok{|\textgreater{}}
  \FunctionTok{filter}\NormalTok{(year }\SpecialCharTok{\textgreater{}=} \DecValTok{2015}\NormalTok{, revenue }\SpecialCharTok{\textgreater{}} \DecValTok{0}\NormalTok{) }\SpecialCharTok{|\textgreater{}}
  \FunctionTok{select}\NormalTok{(year, sector, revenue, employees)}

\CommentTok{\# 3. Komplexe Aggregation mit DuckDB}
\NormalTok{con }\OtherTok{\textless{}{-}} \FunctionTok{dbConnect}\NormalTok{(duckdb}\SpecialCharTok{::}\FunctionTok{duckdb}\NormalTok{())}
\FunctionTok{duckdb\_register\_arrow}\NormalTok{(con, }\StringTok{"btp\_filtered"}\NormalTok{, filtered)}

\NormalTok{result }\OtherTok{\textless{}{-}} \FunctionTok{dbGetQuery}\NormalTok{(con, }\StringTok{"}
\StringTok{  SELECT sector,}
\StringTok{         AVG(revenue) as mean\_revenue,}
\StringTok{         STDDEV(revenue) as sd\_revenue}
\StringTok{  FROM btp\_filtered}
\StringTok{  GROUP BY sector}
\StringTok{  HAVING COUNT(*) \textgreater{}= 30}
\StringTok{"}\NormalTok{)}

\CommentTok{\# 4. Aufräumen}
\FunctionTok{dbDisconnect}\NormalTok{(con, }\AttributeTok{shutdown =} \ConstantTok{TRUE}\NormalTok{)}

\CommentTok{\# 5. Weiterverarbeitung in R}
\NormalTok{result }\OtherTok{\textless{}{-}}\NormalTok{ result }\SpecialCharTok{|\textgreater{}} 
  \FunctionTok{mutate}\NormalTok{(}\AttributeTok{cv =}\NormalTok{ sd\_revenue }\SpecialCharTok{/}\NormalTok{ mean\_revenue)}
\end{Highlighting}
\end{Shaded}

Dieser Workflow kombiniert die Stärken aller Werkzeuge und ist optimal
für FDZ-Bedingungen geeignet.

\bookmarksetup{startatroot}

\chapter{Performanzergebnisse}\label{sec-performance}

\section{Experimentelles Setup}\label{experimentelles-setup}

Die Performanzmessungen wurden auf einem repräsentativen FDZ-System
durchgeführt (2x Intel Xeon Gold 6248R, 48 Kerne, 384 GB RAM, Ubuntu
22.04 LTS). Alle Tests wurden mit R 4.3.2 und aktuellen Versionen der
verwendeten Pakete durchgeführt.

\subsection{Test-Szenarien}\label{test-szenarien}

Die Evaluation umfasst realistische Auswertungsszenarien aus dem BTP:

\begin{enumerate}
\def\labelenumi{\arabic{enumi}.}
\tightlist
\item
  \textbf{Szenario A: Einfache Aggregation} - Jahressummen nach
  Wirtschaftszweig
\item
  \textbf{Szenario B: Komplexe Aggregation} - Multi-dimensionale
  Gruppierung mit mehreren Kennzahlen
\item
  \textbf{Szenario C: Join-Operation} - Verknüpfung zweier großer
  Tabellen
\item
  \textbf{Szenario D: Filterkette} - Mehrere aufeinanderfolgende Filter-
  und Transformationsoperationen
\end{enumerate}

\section{Vergleich der Datenformate}\label{vergleich-der-datenformate}

\subsection{Ladezeiten und
Speicherverbrauch}\label{ladezeiten-und-speicherverbrauch}

\begin{figure}

\centering{

\pandocbounded{\includegraphics[keepaspectratio]{sections/05-performance_files/figure-pdf/fig-format-performance-1.pdf}}

}

\caption{\label{fig-format-performance}Ladezeiten verschiedener
Datenformate (BTP-Datensatz, 5 Mio. Zeilen, 50 Variablen)}

\end{figure}%

\subsection{Aggregations-Performance}\label{aggregations-performance}

\begin{table}

\caption{\label{tbl-aggregation-performance}Ausführungszeiten für
Szenario A: Einfache Aggregation (in Sekunden)}

\centering{

\begin{tabular}{lrrrr}
\toprule
Ansatz & Laden & Aggregation & Gesamt & Speedup\\
\midrule
Base R (read.csv + aggregate) & 45.2 & 18.7 & 63.9 & 1.0\\
data.table (fread + groupby) & 12.1 & 2.3 & 14.4 & 4.4\\
dplyr + CSV & 43.8 & 8.9 & 52.7 & 1.2\\
Arrow + Parquet & 4.8 & 1.2 & 6.0 & 10.7\\
DuckDB + Parquet & 0.3 & 0.8 & 1.1 & 58.1\\
\addlinespace
Arrow + Parquet (partitioniert) & 0.2 & 0.5 & 0.7 & 91.3\\
\bottomrule
\end{tabular}

}

\end{table}%

\section{Skalierungsverhalten}\label{skalierungsverhalten}

\subsection{Performance bei wachsender
Datenmenge}\label{performance-bei-wachsender-datenmenge}

\begin{figure}

\centering{

\pandocbounded{\includegraphics[keepaspectratio]{sections/05-performance_files/figure-pdf/fig-scaling-behavior-1.pdf}}

}

\caption{\label{fig-scaling-behavior}Skalierungsverhalten verschiedener
Ansätze (Szenario B: Komplexe Aggregation)}

\end{figure}%

\section{Join-Performance}\label{join-performance}

\subsection{Verschiedene
Join-Strategien}\label{verschiedene-join-strategien}

Die Verknüpfung großer Datensätze ist eine häufige Operation bei der
Arbeit mit dem BTP:

\begin{figure}

\centering{

\pandocbounded{\includegraphics[keepaspectratio]{sections/05-performance_files/figure-pdf/fig-join-performance-1.pdf}}

}

\caption{\label{fig-join-performance}Join-Performance bei verschiedenen
Datensatzgrößen (Inner Join zweier Tabellen)}

\end{figure}%

\section{Parallelisierungs-Effizienz}\label{parallelisierungs-effizienz}

\subsection{Speedup durch
Multi-Threading}\label{speedup-durch-multi-threading}

\begin{table}

\caption{\label{tbl-parallel-speedup}Speedup durch Parallelisierung
(Szenario D: Filterkette, 10 Mio. Zeilen)}

\centering{

\begin{tabular}{crrrr}
\toprule
Kerne & Arrow (s) & DuckDB (s) & Arrow Speedup & DuckDB Speedup\\
\midrule
1 & 25.3 & 9.8 & 1.0 & 1.0\\
2 & 13.2 & 5.1 & 1.9 & 1.9\\
4 & 7.1 & 2.7 & 3.6 & 3.6\\
8 & 3.9 & 1.5 & 6.5 & 6.5\\
16 & 2.3 & 0.9 & 11.0 & 10.9\\
\addlinespace
32 & 1.8 & 0.7 & 14.1 & 14.0\\
\bottomrule
\end{tabular}

}

\end{table}%

\section{Speichereffizienz}\label{speichereffizienz}

\subsection{Memory Footprint}\label{memory-footprint}

Ein kritischer Faktor in der FDZ-Umgebung ist der Speicherverbrauch:

\begin{figure}

\centering{

\pandocbounded{\includegraphics[keepaspectratio]{sections/05-performance_files/figure-pdf/fig-memory-efficiency-1.pdf}}

}

\caption{\label{fig-memory-efficiency}Peak Memory Usage verschiedener
Ansätze (5 Mio. Zeilen)}

\end{figure}%

\section{Praktische Implikationen}\label{praktische-implikationen}

\subsection{Empfehlungen nach
Anwendungsfall}\label{empfehlungen-nach-anwendungsfall}

Basierend auf den Performanzmessungen lassen sich folgende Empfehlungen
ableiten:

\begin{table}

\caption{\label{tbl-recommendations}Empfohlene Werkzeuge nach
Anwendungsfall}

\centering{

\begin{tabular}{llll}
\toprule
Anwendungsfall & 1. Wahl & 2. Wahl & Begründung\\
\midrule
Kleine Datensätze (<1M Zeilen) & data.table & Base R & Einfachheit, gute Performance\\
Explorative Analyse, häufige Iterationen & DuckDB & Arrow & SQL-Komfort, schnelle Iteration\\
Große Aggregationen (>5M Zeilen) & DuckDB & Arrow + Parquet & Beste Skalierung, niedrigster Speicherverbrauch\\
Komplexe Joins & DuckDB & data.table & Optimierter Hash-Join, Index-Support\\
Produktions-Pipelines & Arrow + Parquet & DuckDB & Persistenz, Partitionierung, Reproduzierbarkeit\\
\addlinespace
Speicher-limitierte Umgebung & DuckDB & Arrow (streaming) & Out-of-core Verarbeitung möglich\\
\bottomrule
\end{tabular}

}

\end{table}%

\section{Zusammenfassung der
Erkenntnisse}\label{zusammenfassung-der-erkenntnisse}

Die empirischen Messungen zeigen deutlich:

\begin{enumerate}
\def\labelenumi{\arabic{enumi}.}
\item
  \textbf{Moderne Formate dominieren:} Arrow/Parquet und DuckDB bieten
  um den Faktor 10-90x bessere Performance als traditionelle Ansätze
\item
  \textbf{Partitionierung lohnt sich:} Bei großen Datensätzen reduziert
  Partitionierung die Zugriffszeiten erheblich
\item
  \textbf{Query Optimization wichtig:} Lazy Evaluation und Predicate
  Pushdown ermöglichen dramatische Speedups
\item
  \textbf{Parallelisierung skaliert gut:} Bis zu 14x Speedup mit 32
  Kernen möglich
\item
  \textbf{Speicher ist der Engpass:} Moderne Ansätze benötigen 3-7x
  weniger RAM als Base R
\end{enumerate}

Diese Ergebnisse zeigen, dass durch den Einsatz moderner
Auswertungsmethoden die Effizienz der Datenanalyse in
Forschungsdatenzentren erheblich gesteigert werden kann.

\bookmarksetup{startatroot}

\chapter{Ausblick}\label{sec-ausblick}

\section{Zusammenfassung der
Erkenntnisse}\label{zusammenfassung-der-erkenntnisse-1}

Die vorliegende Untersuchung hat gezeigt, dass moderne
spaltenorientierte Technologien erhebliche Vorteile für die Analyse
großer Datensätze in der amtlichen Statistik bieten. Die Kombination von
\textbf{Apache Arrow}, \textbf{DuckDB} und \textbf{Parquet} ermöglicht:

\begin{itemize}
\tightlist
\item
  \textbf{Drastische Reduktion der Speicheranforderungen} (typisch 5-10x
  Kompression gegenüber CSV)
\item
  \textbf{Signifikante Beschleunigung der Analysen} durch
  spaltenorientierte Verarbeitung
\item
  \textbf{Out-of-Core-Verarbeitung} für Datensätze, die den verfügbaren
  RAM übersteigen
\item
  \textbf{Effiziente Query-Optimierung} durch Lazy Evaluation und
  Predicate Pushdown
\end{itemize}

\section{Implikationen für die
Forschungsdatenzentren}\label{implikationen-fuxfcr-die-forschungsdatenzentren}

\subsection{Infrastrukturentwicklung}\label{infrastrukturentwicklung}

Die FDZ sollten die Integration dieser Technologien vorantreiben:

\begin{enumerate}
\def\labelenumi{\arabic{enumi}.}
\tightlist
\item
  \textbf{Datenbereitstellung}: Zusätzlich zu SAS/Stata-Formaten auch
  Parquet-Versionen der Datensätze anbieten
\item
  \textbf{Software-Ausstattung}: R mit Arrow- und DuckDB-Paketen auf
  allen Systemen verfügbar machen
\item
  \textbf{Schulungsangebote}: Workshops und Dokumentation für Forschende
  bereitstellen
\item
  \textbf{Best-Practice-Katalog}: Erfolgreiche Analysemuster
  dokumentieren und teilen
\end{enumerate}

\subsection{Nutzerunterstützung}\label{nutzerunterstuxfctzung}

Forschende profitieren von:

\begin{itemize}
\tightlist
\item
  \textbf{Kürzeren Analysezeiten}: Mehr explorative Analyse in kürzerer
  Zeit möglich
\item
  \textbf{Größere Datensätze handhabbar}: Vollsatz-Analysen statt
  Stichproben
\item
  \textbf{Geringere Ressourcenkonflikte}: Effizientere Nutzung geteilter
  Infrastruktur
\item
  \textbf{Reproduzierbare Workflows}: Parquet-basierte Pipelines gut
  dokumentierbar
\end{itemize}

\section{Technologische
Entwicklungen}\label{technologische-entwicklungen}

\subsection{Kurz- bis mittelfristig}\label{kurz--bis-mittelfristig}

\textbf{Weitere Optimierungen} der bestehenden Tools sind zu erwarten:

\begin{itemize}
\tightlist
\item
  \textbf{Arrow-Flight}: Effizienter Datentransfer zwischen Prozessen
  und Systemen
\item
  \textbf{Substrait}: Standardisierte Query-Repräsentation für
  Tool-übergreifende Optimierung
\item
  \textbf{DuckDB Extensions}: Spezielle Module für statistische Analysen
\item
  \textbf{GPU-Beschleunigung}: Erste Implementierungen für
  Arrow-basierte Berechnungen
\end{itemize}

\subsection{Langfristig}\label{langfristig}

\textbf{Neue Paradigmen} könnten zusätzliche Verbesserungen bringen:

\begin{itemize}
\tightlist
\item
  \textbf{Föderated Queries}: Analyse über verteilte Datensätze ohne
  zentrale Zusammenführung
\item
  \textbf{Adaptive Execution}: Automatische Wahl der optimalen
  Verarbeitungsstrategie
\item
  \textbf{Streaming Analytics}: Echtzeitverarbeitung großer Datenströme
\item
  \textbf{Privacy-Preserving Computation}: Analysen auf verschlüsselten
  Daten
\end{itemize}

\section{Offene Forschungsfragen}\label{offene-forschungsfragen}

Weitere Untersuchungen sollten folgende Aspekte betrachten:

\begin{enumerate}
\def\labelenumi{\arabic{enumi}.}
\item
  \textbf{Statistische Software-Integration}: Wie können
  R-Modellierungsfunktionen optimal mit Arrow/DuckDB kombiniert werden?
\item
  \textbf{Datenschutz-Aspekte}: Welche Implikationen haben neue
  Datenformate für die Output-Kontrolle im FDZ?
\item
  \textbf{Qualitätssicherung}: Wie können Korrektheit und
  Reproduzierbarkeit bei komplexen Arrow/DuckDB-Workflows sichergestellt
  werden?
\item
  \textbf{Heterogene Systeme}: Wie lassen sich verschiedene Datenquellen
  und -formate optimal integrieren?
\end{enumerate}

\section{Handlungsempfehlungen}\label{handlungsempfehlungen}

\subsection{Für
Forschungsdatenzentren}\label{fuxfcr-forschungsdatenzentren}

\textbf{Pilotprojekte} sollten initiiert werden:

\begin{itemize}
\tightlist
\item
  Ausgewählte häufig genutzte Datensätze (wie BTP) in Parquet
  bereitstellen
\item
  Erfolgsmetriken definieren (Analysezeit, Ressourcenverbrauch,
  Nutzerzufriedenheit)
\item
  Feedback von Power-Usern systematisch einholen
\item
  Erkenntnisse in Qualitätsstandards überführen
\end{itemize}

\subsection{Für Forschende}\label{fuxfcr-forschende}

\textbf{Schrittweise Einführung} wird empfohlen:

\begin{enumerate}
\def\labelenumi{\arabic{enumi}.}
\tightlist
\item
  Kleine Analyseprojekte als Lernumgebung nutzen
\item
  Kritische Analysen zunächst parallel mit bewährten Methoden
  durchführen
\item
  Performanzgewinne und Limitationen dokumentieren
\item
  Bei Erfolg schrittweise auf größere Projekte übertragen
\end{enumerate}

\section{Schlussbetrachtung}\label{schlussbetrachtung}

Die Analyse großer Datensätze in Forschungsdatenzen tren steht vor einem
Technologiewandel. Die hier untersuchten Werkzeuge -- \textbf{Arrow,
DuckDB und Parquet} -- bieten das Potenzial, die Effizienz erheblich zu
steigern und neue Analysemöglichkeiten zu eröffnen.

Der Übergang erfordert jedoch eine \textbf{koordinierte Anstrengung} von
Infrastrukturbetreibern, Software-Entwicklern und Forschenden. Die
systematische Evaluation, Dokumentation und Verbreitung von Best
Practices wird entscheidend für den Erfolg sein.

Die amtliche Statistik kann von diesen Entwicklungen in mehrfacher
Hinsicht profitieren: \textbf{schnellere Erkenntnisgewinnung},
\textbf{umfassendere Analysen} und \textbf{effizientere
Ressourcennutzung}. Dies trägt letztlich zur Stärkung der
evidenzbasierten Politik- und Gesellschaftsberatung bei.


\backmatter
\printbibliography



\end{document}
